\قسمت{پراکسی‌ها}

فرض کنید طراحی یک سیستم به شما سپرده شده است، این سیستم چند نمونه از یک سرویس در حال اجرا دارد که شما می‌بایست تقاضاها را بین این نمونه‌ها تقسیم کنید. برای اینکار از چه نوع پراکسی استفاده می‌کنید؟
در رابطه با تفاوت \متن‌لاتین{forward proxy} و \متن‌لاتین{reverse proxy} توضیح دهید.

\نمره{۳}

% پاسخ پیشنهادی است و پیش از استفاده در تصحیح نیاز به بازبینی دارد.
\begin{پاسخ}

برای تقسیم تقاضاها میان چند نمونه از یک سرویس، \متن‌لاتین{reverse proxy} استفاده می‌شود که در این کاربرد همان توزیع‌کننده‌ی بار\پانویس{Load Balancer} است.

\شروع{فقرات}
\فقره \متن‌سیاه{\متن‌لاتین{forward proxy}}: جلوی کلاینت‌ها می‌ایستد، کلاینت آن را می‌شناسد و تنظیمش می‌کند، و هویت کلاینت را از سرور پنهان می‌کند.
\فقره \متن‌سیاه{\متن‌لاتین{reverse proxy}}: جلوی سرورها می‌ایستد، کلاینت آن را خودِ سرور می‌پندارد، و ساختار پشت آن را پنهان می‌کند.
\پایان{فقرات}

\شروع{فقرات}
\فقره انتخاب \متن‌لاتین{reverse proxy}: ۵۰ درصد نمره
\فقره تفاوت دو پراکسی از نظر اینکه جلوی چه کسی قرار می‌گیرند: ۵۰ درصد نمره
\پایان{فقرات}

\end{پاسخ}
